% Concise, code-aligned details for the LPB implementation of DAPRO.
% This section is intended to follow the main-text method description.

\subsection{Implementation details for LPB \ttmethod}
\label{sec:lpb_dapro_implementation}

We give only the details needed to reproduce the LPB allocator that are not
specified in Section~\ref{sec:method}. Recall that
$b_i:=\min\{T_i,\qprior(X_i)\}$. After fully observing \(\calIone\), the
policy-shape budget is
\(
\Bbar=(B-\sum_{i\in\calIone}b_i)/(|\calIcrc|+|\calItwo|)
\), as in the main text. The rows in \(\calIcrc\) are also fully observed,
but their realized costs enter the CRC selector in
Eq.~\eqref{eq:dapro_main_crc_selector}; subtracting them again in \(\Bbar\)
would count the same cost twice.

\paragraph{Soft-prefix target and score.}
At an observed prefix \(H_i(t)\), the survival model supplies the causal
current-event probability
\begin{equation}
    \widehat h_i(t)
    :=\widehat{\mathbb P}\!\left(T_i=t\mid T_i\geq t,H_i(t)\right),
    \qquad S_i(t):=\widehat h_i(t).
    \label{eq:lpb_impl_hazard_score}
\end{equation}
The numerator of Eq.~\eqref{eq:opt_prob} is implemented with the soft-prefix
coefficient
\begin{equation}
    a_i(t)
    :=\mathbb I\{t\leq b_i\}\widehat h_i(t)
      \frac{\mathbb I\{t<\fhat_\alpha(X_i)\}+\delta}{1+\delta}.
    \label{eq:lpb_impl_soft_mass}
\end{equation}
The leading term estimates the probability that the event occurs at turn
\(t\) and contributes to LPB miscoverage; the small global term \(\delta\)
prevents all other observable events from receiving zero optimization signal.
Both \(S_i(t)\) and \(a_i(t)\) use only the history available before the
continuation decision at turn \(t\).

\paragraph{Deployable score map and numerical optimization.}
At each turn, the active Phase-I scores are divided at their empirical median.
The optimizer assigns one continuation probability to each of the two bins,
constrained to be nondecreasing in the score. Ties enter the upper bin, an
empty bin inherits the nearest nonempty value, and a turn with no active
Phase-I row is continued deterministically.

Equation~\eqref{eq:opt_prob} is solved in log conditional probabilities. For
each dual multiplier, at most \(30\) coordinate sweeps solve the monotone
subproblems by pool-adjacent-violators updates, stopping when the largest
log-probability change is below \(5\times10^{-4}\). An outer geometric
bisection uses at most \(60\) iterations and stops when the mean-cost error is
at most \(10^{-9}\max\{\Bbar,1\}\).

Projection to two score bins can perturb the Phase-I cost. The implementation
therefore applies one shared logit shift to the projected cumulative reaches.
If \(\rho_i^{\rm raw}(t):=\prod_{s=1}^tP_i^{\rm raw}(s)\), define
\begin{equation}
    \rho_i^{\rm base}(t)
    :=\varepsilon+(1-\varepsilon)
      \min_{u\leq t}
      \sigma\!\left(\operatorname{logit}\rho_i^{\rm raw}(u)+\kappa\right).
    \label{eq:lpb_impl_cumulative_correction}
\end{equation}
Here \(\kappa\) is the largest shift whose empirical cost on \(\calIone\) is
at most \(\Bbar\); it is found by at most \(80\) bisection steps, with shift
tolerance \(10^{-10}\) and cost tolerance \(10^{-7}\). We recover
\(P_i^{\rm base}(t)=\rho_i^{\rm base}(t)/\rho_i^{\rm base}(t-1)\), where
\(\rho_i^{\rm base}(0)=1\). The cumulative minimum makes the reach path
temporally coherent, and \(\varepsilon>0\) gives positive inclusion
probabilities.

\paragraph{Independent CRC calibration.}
CRC requires every candidate policy to have bounded per-row expected cost.
Using only \(\calIone\), the implementation first pools the soft masses as
\(\bar a(t):=N_1^{-1}\sum_{i\in\calIone}a_i(t)\) and fits a shared
nonincreasing cumulative-reach envelope \(e(1)\geq\cdots\geq e(\maxl)\) by
\begin{equation}
    \min_e\ \sum_{t=1}^{\maxl}\frac{\bar a(t)}{e(t)}
    \quad\text{s.t.}\quad
    1\geq e(1)\geq\cdots\geq e(\maxl)\geq\varepsilon,
    \qquad \sum_{t=1}^{\maxl}e(t)\leq C_{\max}.
    \label{eq:lpb_impl_shared_envelope}
\end{equation}
This one-dimensional problem is solved by antitonic PAV and a common scale.
The resulting \(e\) is frozen before examining \(\calIcrc\). For any new
score history, let
\begin{equation}
    \rho_i^{\rm cap}(t)
      :=\min\{\rho_i^{\rm base}(t),e(t)\},
    \qquad
    P_i^{\rm cap}(t)
      :=\frac{\rho_i^{\rm cap}(t)}{\rho_i^{\rm cap}(t-1)},
    \qquad \rho_i^{\rm cap}(0):=1.
    \label{eq:lpb_impl_capped_policy}
\end{equation}
Equivalently, the capped policy is the causal, stateful map
\(M_t^{\rm cap}(S_i(1{:}t)):=P_i^{\rm cap}(t)\). It remains deployable
because it needs only the current base probability and the previous cumulative
reach. Moreover,
\(
\bfunc(M^{\rm cap}(S_i))
=\sum_{t\leq\qprior(X_i)}\rho_i^{\rm cap}(t)
\leq\sum_t e(t)\leq C_{\max}
\)
for every trajectory, which supplies the bounded loss required by CRC.

The nested CRC family contracts this capped cumulative reach toward the
positivity floor:
\begin{equation}
    \rho_i(t;\lambda)
    :=\varepsilon+\lambda
      \bigl(\rho_i^{\rm cap}(t)-\varepsilon\bigr),
    \qquad
    M_t(S_i(1{:}t);\lambda)
    :=P_i(t;\lambda)
    :=\frac{\rho_i(t;\lambda)}{\rho_i(t-1;\lambda)},
    \quad \lambda\in[0,1].
    \label{eq:lpb_impl_crc_family}
\end{equation}
Thus increasing \(\lambda\) increases every cumulative reach and hence
\(\bfunc(M(S_i;\lambda))\), while retaining the allocation learned in
Phase~I. The implementation evaluates this frozen family on \(\calIcrc\)
and selects \(\widehat\lambda\) by
Eq.~\eqref{eq:dapro_main_crc_selector}. Without CRC, it deploys
\(M^{\rm base}\) directly, with no projection-error reserve; this version
satisfies the empirical Phase-I constraint but does not claim the CRC
finite-sample budget guarantee.

\paragraph{Bookkeeping, defaults, and cost.}
Once the event or \(\qprior(X_i)\) is reached, the stored censoring time is
\(C_i=\qprior(X_i)\), although an earlier event consumes only \(T_i\)
interactions. For \(i\in\calIone\cup\calIcrc\), the implementation sets
\(C_i=\qprior(X_i)\) and \(w(i)=1\). The reported LPB experiments use two
score bins, \(\delta=10^{-3}\), \(\varepsilon=0.005\), and zero
projection-error margin. The CRC version uses
\(C_{\max}=\min\{\maxl,2B/|\calI|\}\) and \(401\) candidate values of
\(\lambda\); these choices vary only in their designated ablations.

For completeness, let \(H=\maxl\), \(n_1=|\calIone|\),
\(n_c=|\calIcrc|\), \(n_2=|\calItwo|\), and let \(J\) be the number of CRC
candidates. With the fixed iteration limits above, the worst-case running
time is
\begin{equation}
    \mathcal O\!\left(
      60\cdot30\,n_1H^2+n_1H\log n_1+Jn_cH+n_2H
    \right),
\end{equation}
where the first term is the coordinate solver, the second fits the score maps,
and the last two evaluate CRC and deploy the policy. Memory usage is
\(\mathcal O((n_1+n_c+n_2)H+Jn_c)\). These are conservative worst-case
bounds; early stopping and convergence usually reduce both computation and
acquisition substantially.
